% Chapter 6 --- The tensor product and the symmetric monoidal closed
% structure (paper §5.3–§5.5, iteration 2).
%
% Mirrors:
%   theories/homs/icones_iso.v       isos in ICones (groupoid structure)
%   theories/mcones/test_pullback.v  Def 3.13 test-side characterisation
%   theories/homs/linhom_functor.v   Prop 5.8, the internal-hom functor
%   theories/axioms/saft_interface.v the ~18 staged SAFT/tensor Parameters
%   theories/homs/tensor.v           §5.4: τ, x⊗y, Eq 5.1, Thm 5.12/5.13,
%                                    Prop 5.14
%   theories/homs/smcc.v             §5.5: α/λ/ρ/σ + derived coherence,
%                                    Thm 5.15 (ICones is SMCC)
%
% Honesty note: tensor.v and smcc.v are mechanised *modulo* the staging
% interface saft_interface.v (the SAFT/Yoneda representability content).
% The structural-iso EXISTENCE is staged (\notready); the monoidal
% coherence and the consequences are genuinely PROVED from the interface
% (\rocqok). The MVP headline (Thm 6.5, Chapter \ref{ch:skern}) does not
% depend on this interface.

%%%%%%%%%%%%%%%%%%%%%%%%%%%%%%%%%%%%%%%%%%%%%%%%%%%%%%%%%%%%%%%%%%%%%%%%
\chapter{The tensor product and the symmetric monoidal closed
structure}\label{ch:tensor}

This chapter formalises paper \S5.3--\S5.5: the symmetric monoidal
closed (SMC closed) structure of $\ICones$, with the tensor product
$\otimes$ as its monoidal product and the internal hom $\lin$ of
Chapter~\ref{ch:linhom} as its closed structure (paper
Theorem~5.15). Paper \S5 obtains the cartesian/monoidal closed
structure of $\ICones$ in one of two ways: \emph{either} concretely
--- as for the internal hom $C \lin D$, built by hand in
Chapter~\ref{ch:linhom} --- \emph{or} abstractly, via Freyd's Special
Adjoint Functor Theorem (SAFT) applied to the limit-preserving functor
$C \lin -\;: \ICones \to \ICones$. The paper reserves the abstract,
SAFT-based argument for the tensor product $\otimes$ (and, beyond the
scope of this chapter, for the exponential $\mathord{!}$). This
chapter formalises the SMC closed structure that results once
$\otimes$ is available.

The material lives under \texttt{theories/homs/icones\_iso.v} (isos in
$\ICones$), \texttt{theories/mcones/test\_pullback.v} (the Def~3.13
test-side characterisation), \texttt{theories/homs/linhom\_functor.v}
(the morphism-level internal-hom functor, Prop~5.8),
\texttt{theories/homs/tensor.v} (\S5.4) and
\texttt{theories/homs/smcc.v} (\S5.5), with the SAFT/tensor universal
property collected as a staging interface in
\texttt{theories/axioms/saft\_interface.v}.

\section{Staging strategy: the SAFT/tensor interface}%
\label{sec:tensor-staging}

This is the central design decision of the chapter, and we state it up
front. The \emph{representability} content that the paper extracts
from SAFT --- the existence of a tensor object $B \otimes C$
representing the functor $D \mapsto \ICones(B, C \lin D)$, and the
existence of the structural isomorphisms $\alpha/\lambda/\rho/\sigma$
produced by Yoneda (paper Lemma~5.5) --- is the deep, non-elementary
input to the monoidal structure. The Rocq development \emph{stages}
this content: it is collected as a small interface of approximately
$18$ \texttt{Parameter}s in
\texttt{theories/axioms/saft\_interface.v}, and the \emph{consequences}
(\texttt{tensor.v}, \texttt{smcc.v}) are mechanised \emph{modulo} this
interface. Every \texttt{Parameter} in the interface is an
intentional, temporary axiom, flagged in the source with
\texttt{STAGING: discharge via M-SAFT, PLAN \S13.1; then delete}. No
other part of the project is left unproved, and there is no
\Admitted{} anywhere.

We are precise about what is staged and what is proved:
\begin{enumerate}
  \item[(a)] The substantive monoidal coherence --- the triangle,
    pentagon and braiding-hexagon identities, and the symmetry law
    $\sigma \circ \sigma = \id$ --- is \emph{genuinely proved} from
    the interface, via the pure-tensor tree-extensionality of paper
    Proposition~5.14 (\S\ref{sec:tensor-ext}).
  \item[(b)] Only the structural-iso \emph{existence} (the
    Lemma-5.5/Yoneda content) and the universal property of $\otimes$
    are staged: the tensor object former, the hom-bijection $\Phi$ and
    its naturality, the Thm~5.12 element-level iso, the Thm~5.13 norm
    equation, and the four structural isos $\alpha/\lambda/\rho/\sigma$
    of Eqs~5.2--5.4 together with their pure-tensor computation laws.
  \item[(c)] The plan (\texttt{PLAN.md} \S13.1, milestone
    ``M-SAFT'') is to discharge every \texttt{Parameter} by proving
    SAFT for $\ICones$, after which each \texttt{Parameter} is
    replaced by a theorem and \texttt{saft\_interface.v} is deleted,
    reaching zero project axioms.
  \item[(d)] The MVP headline theorem --- paper Theorem~6.5, the fully
    faithful embedding $\Skern \to \ICones$ of
    Chapter~\ref{ch:skern} --- does \emph{not} depend on this
    interface; the staging is confined to the $\otimes$/SMC-closed
    development of this chapter.
\end{enumerate}

\paragraph{The $18$ staged \texttt{Parameter}s, grouped.}
The interface is stated in a single section over a real type $R$ and a
measurable subcategory $\Ar$, with local notations $C \lin D :=
\texttt{linhom\_car}\,\Ar\,C\,D$ and $B \otimes C := \texttt{tensor}\,\Ar\,B\,C$.
The \texttt{Parameter}s fall into five groups.

\begin{itemize}
  \item \emph{The representing object.}\label{def:tensor-object}
    \rocq{Icones.axioms.saft\_interface.tensor}\notready{} ---
    the opaque tensor object former $B \otimes C$ (paper \S5.4,
    Remark~5.1, which gives no explicit carrier).
  \item \emph{The hom-bijection $\Phi$ and its inverse (paper Thm~5.9 /
    Eq~5.1).} The forward map
    \rocq{Icones.axioms.saft\_interface.tensor\_curry}\notready{}
    $:\ICones(B \otimes C, D) \to \ICones(B, C \lin D)$, the inverse
    \rocq{Icones.axioms.saft\_interface.tensor\_uncurry}\notready{},
    and the two round-trip cancellations
    \rocq{Icones.axioms.saft\_interface.tensor\_curryK}\notready{},
    \rocq{Icones.axioms.saft\_interface.tensor\_uncurryK}\notready{}.
  \item \emph{Naturality of $\Phi$.} Naturality in $D$ stated
    existentially
    (\rocq{Icones.axioms.saft\_interface.tensor\_curry\_natural\_D}\notready{}),
    naturality in $B$
    (\rocq{Icones.axioms.saft\_interface.tensor\_curry\_natural\_B}\notready{}),
    and the \emph{concrete} naturality in $D$ whose existential witness
    is the explicit hom-functor action $C \lin h$
    (\rocq{Icones.axioms.saft\_interface.tensor\_curry\_natural\_post}\notready{}).
    The last one is the workhorse: it yields Eq~5.1 below.
  \item \emph{Element-level packaging and the norm equation.} The
    Thm~5.12 iso of integrable cones
    \rocq{Icones.axioms.saft\_interface.tensor\_hom\_iso}\notready{}
    and the Thm~5.13 norm equation
    \rocq{Icones.axioms.saft\_interface.tensor\_normM}\notready{}.
  \item \emph{Structural-iso existence and their computation laws (paper
    Eqs~5.2--5.4).} The associator
    \rocq{Icones.axioms.saft\_interface.tensor\_assoc\_iso}\notready{}
    with
    \rocq{Icones.axioms.saft\_interface.tensor\_assocE}\notready{}, the
    unitors
    \rocq{Icones.axioms.saft\_interface.tensor\_lunit\_iso}\notready{}
    /
    \rocq{Icones.axioms.saft\_interface.tensor\_runit\_iso}\notready{}
    with
    \rocq{Icones.axioms.saft\_interface.tensor\_lunitE}\notready{} /
    \rocq{Icones.axioms.saft\_interface.tensor\_runitE}\notready{}, and
    the braiding
    \rocq{Icones.axioms.saft\_interface.tensor\_braid\_iso}\notready{}
    with
    \rocq{Icones.axioms.saft\_interface.tensor\_braidE}\notready{}.
\end{itemize}

\begin{remark}[A known gap in the interface, deferred to M-SAFT]
Naturality of $\Phi$ in the $C$ slot is \emph{not} stated, because it
needs the contravariant functorial action $\lin$ in its left argument,
$\texttt{linhom\_pre\_icones}\,v : (C \lin D) \to (C' \lin D)$, in a
shape the interface does not pin down. The tensor milestone of this
chapter consumes only naturality in $D$ and $B$; naturality in $C$ is
needed only for the symmetry coherence, which is instead derived from
the pure-tensor computation law $\sigma(x \otimes y) = y \otimes x$.
The $C$-slot square is to be added when SAFT lands (a \texttt{TODO} in
\texttt{saft\_interface.v}).
\end{remark}

\section{Prerequisites for $\otimes$ (paper \S5.3)}%
\label{sec:tensor-prereq}

Three pieces of infrastructure feed the tensor development. All three
are genuinely proved.

\subsection{Isomorphisms in $\ICones$}\label{sec:icones-iso}

The structural equivalences of the monoidal structure are stated as
isos in the concrete category $\ICones$ (Chapter~\ref{ch:icones}),
whose objects are $\ICone$s, whose morphisms are
$\icones_{\mathrm{hom}}$, with identity \texttt{icones\_id},
composition \texttt{icones\_comp} (where $\texttt{icones\_comp}\,g\,f$
is ``$g$ after $f$''), and equality of morphisms reduced to pointwise
equality by \texttt{icones\_hom\_eq}.

\begin{definition}[Iso of integrable cones]\label{def:icones-iso}
\rocq{Icones.homs.icones\_iso.icones\_iso}\rocqok\uses{sec:icones}
An isomorphism $B \cong C$ in $\ICones$ is a record
\rocq{Icones.homs.icones\_iso.icones\_iso} packaging a forward
morphism $\texttt{iso\_fwd} : \icones_{\mathrm{hom}}\,\Ar\,B\,C$, a
backward morphism $\texttt{iso\_bwd} : \icones_{\mathrm{hom}}\,\Ar\,C\,B$,
and the two round-trip equations
$\texttt{icones\_comp}\,(\texttt{iso\_bwd})\,(\texttt{iso\_fwd}) =
\texttt{icones\_id}\,\Ar\,B$ (\texttt{iso\_fwdK}) and
$\texttt{icones\_comp}\,(\texttt{iso\_fwd})\,(\texttt{iso\_bwd}) =
\texttt{icones\_id}\,\Ar\,C$ (\texttt{iso\_bwdK}).
\end{definition}

The smart constructor
\rocq{Icones.homs.icones\_iso.icones\_isoP} builds an iso from the two
cancellation equations stated at the level of \texttt{icones\_comp};
the Yoneda-style companion
\rocq{Icones.homs.icones\_iso.icones\_iso\_of\_cancel} builds it from
\emph{pointwise} cancellation, upgrading via \texttt{icones\_hom\_eq}.

\begin{lemma}[Groupoid structure and bijectivity]\label{lem:icones-iso-groupoid}
\rocq{Icones.homs.icones\_iso.icones\_iso\_refl}\rocqok
\rocq{Icones.homs.icones\_iso.icones\_iso\_sym}\rocqok
\rocq{Icones.homs.icones\_iso.icones\_iso\_trans}\rocqok
\uses{def:icones-iso}
The isos form a groupoid: the identity iso $B \cong B$
(\rocq{Icones.homs.icones\_iso.icones\_iso\_refl}), symmetry $B \cong C
\Rightarrow C \cong B$ by swapping forward and backward
(\rocq{Icones.homs.icones\_iso.icones\_iso\_sym}), and transitivity by
composition (\rocq{Icones.homs.icones\_iso.icones\_iso\_trans}).
Moreover the underlying forward function is bijective: it has the
backward function as a two-sided pointwise inverse
(\rocq{Icones.homs.icones\_iso.iso\_can},
\rocq{Icones.homs.icones\_iso.iso\_can\'}), hence is injective
(\rocq{Icones.homs.icones\_iso.iso\_fwd\_inj}) and bijective
(\rocq{Icones.homs.icones\_iso.iso\_fwd\_bij}).
\end{lemma}

\subsection{The test-pullback characterisation (paper Def 3.13)}%
\label{sec:test-pullback}

The morphism-level internal-hom functor of \S5.3 needs the
\emph{test-side} reformulation of an $\mcones_{\mathrm{hom}}$. Paper
Definition~3.13 (the ``equivalently'' bullet) says a
$\cones_{\mathrm{hom}}\;g : B \to C$ is in $\MCones$ iff, for every
test $m \in \Mtest^C_Z$ and every measurable path $\varphi : W \to B$,
the function $(s, r) \mapsto m\bigl(s, g(\varphi(r))\bigr)$ is
measurable on $Z \times W$. We package the forward direction in two
shapes.

\begin{lemma}[Test-pullback along an $\mcones_{\mathrm{hom}}$, paper
Def 3.13]\label{lem:test-pullback-meas}
\rocq{Icones.mcones.test\_pullback.test\_pullback\_meas}\rocqok
\rocq{Icones.mcones.test\_pullback.test\_pullback\_meas\_bivar}\rocqok
\uses{def:mcone}
Let $g : \mcones_{\mathrm{hom}}\,\Ar\,B\,C$. For a target test $m \in
\Mtest^C_Z$ and a measurable path $\varphi : W \to B$, the function
$(s, r) \mapsto m\bigl(s, g(\varphi(r))\bigr)$ is measurable
(\rocq{Icones.mcones.test\_pullback.test\_pullback\_meas}). The
bivariate variant
(\rocq{Icones.mcones.test\_pullback.test\_pullback\_meas\_bivar})
folds the path argument out of a pair $(s, r) : Z \times W$ through
$\texttt{ar\_prod}$, giving the exact shape arising in \S5.3's
$\varphi(s, r) = m\bigl(s, g(\delta(s, r))\bigr)$ for a bivariate path
$\delta$.
\end{lemma}

\begin{proof}[Proof sketch]
Both reduce to \texttt{measurable\_test\_path\_section} of M3 applied
to the post-composed path $g \circ \varphi$, whose measurability is
exactly $\texttt{mcones\_hom\_pres\_path}\,g$
(\rocq{Icones.mcones.test\_pullback.test\_pullback\_path}). The
bivariate version reindexes $m$ along $\texttt{ar\_prod\_fst}$ to a
test at the product arity $\texttt{ar\_prod}\,Z\,W$, applies the basic
section, then folds the input pair back via the
$\texttt{ar\_prod\_cast}$ / \texttt{test\_reindex} pattern shared with
\texttt{bilin.v} and \texttt{linhom.v}.
\end{proof}

\begin{lemma}[Test-pullback continuity, paper \S5.3 / Def 3.13]%
\label{lem:test-pullback-cont}
\rocq{Icones.mcones.test\_pullback.test\_pullback\_cont}\rocqok
\uses{lem:test-pullback-meas}
For a $\cones_{\mathrm{hom}}\;g : D_1 \to D_2$, a test $m \in
\Mtest^{D_2}_Z$, and a $\preceq$-increasing unit-ball chain $(a_n)_n$
in $D_1$,
\[
  m\bigl(s, g(\textstyle\bigsup_n a_n)\bigr)
  = \sup_n\, m\bigl(s, g(a_n)\bigr).
\]
\end{lemma}

\begin{proof}[Proof sketch]
The $g$-image of a unit-ball increasing chain is again a unit-ball
increasing chain ($g$ is linear hence increasing,
\rocq{Icones.mcones.test\_pullback.cones\_hom\_image\_chain}, and
norm-decreasing,
\rocq{Icones.mcones.test\_pullback.cones\_hom\_image\_ub1}); push the
cone-sup through $g$ by $\omega$-continuity of $g$, then apply the
``test-of-sup'' identity
\rocq{Icones.mcones.test\_pullback.test\_of\_sup} (a test is monotone
and $\omega$-continuous, so its value at a $\texttt{cone\_sup\_ball}$
is the real supremum of the chain of test values).
\end{proof}

\subsection{The internal-hom functor on morphisms (paper Prop 5.8)}%
\label{sec:linhom-functor}

\S5.3 upgrades the internal hom $\lin$ from an object map to a functor
of \emph{two} variables (contravariant in the source, covariant in the
target). Given $h : C_2 \to C_1$ and $g : D_1 \to D_2$, the object map
$\Lambda := \texttt{linhom\_map\_fun}\,h\,g$ sends $f \mapsto g \circ f
\circ h$, so $(\Lambda f)\,x = g\bigl(f(h\,x)\bigr)$.

\begin{lemma}[Paper Prop 5.8, the hom-functor action]%
\label{lem:linhom-map-icones}
\rocq{Icones.homs.linhom\_functor.linhom\_map\_icones}\rocqok
\uses{def:linhom-car, lem:test-pullback-meas, lem:test-pullback-cont}
The map $\Lambda = (f \mapsto g \circ f \circ h)$ packages as a
morphism
$\texttt{linhom\_map\_icones}\,h\,g :
\icones_{\mathrm{hom}}\,\Ar\,(C_1 \lin D_1)\,(C_2 \lin D_2)$, with
$\bigl(\texttt{linhom\_map\_icones}\,h\,g\bigr)\,f = g \circ f \circ h$
(\rocq{Icones.homs.linhom\_functor.linhom\_map\_iconesE}).
\end{lemma}

\begin{proof}[Proof sketch]
The construction climbs the three layers $\Cone \to \MCone \to \ICone$
on $\Lambda$:
\begin{itemize}
  \item \emph{$\cones_{\mathrm{hom}}$}
    (\rocq{Icones.homs.linhom\_functor.linhom\_map\_cones}): linearity
    (\rocq{Icones.homs.linhom\_functor.linhom\_map\_linear}),
    norm-decrease $\cnorm{\Lambda f} \le \cnorm{f}$ from the chain
    $\cnorm{g(f(h\,x))} \le \cnorm{f(h\,x)} \le \cnorm{f}\,\cnorm{x}$
    (\rocq{Icones.homs.linhom\_functor.linhom\_map\_norm\_le1}), and
    $\omega$-continuity
    (\rocq{Icones.homs.linhom\_functor.linhom\_map\_continuous}) ---
    which, by test-separation (Mssep) of $D_2$, reduces to pushing the
    cone-sup through $g$ via
    \autoref{lem:test-pullback-cont}.
  \item \emph{$\mcones_{\mathrm{hom}}$}
    (\rocq{Icones.homs.linhom\_functor.linhom\_map\_mcones}): the crux,
    measurable-path preservation \emph{in $f$}
    (\rocq{Icones.homs.linhom\_functor.linhom\_map\_pres\_path}). For a
    measurable path $\Phi$ of $C_1 \lin D_1$, the bivariate $D_1$-path
    $\delta(s, r) := (\Phi r)\bigl(h(\gamma_2 s)\bigr)$ is measurable
    (by testing $\Phi$ against the linhom test of $h \circ \gamma_2$),
    and pulling a $D_2$-test back along $g$ keeps it measurable by the
    bivariate
    \autoref{lem:test-pullback-meas}
    (\rocq{Icones.mcones.test\_pullback.test\_pullback\_meas\_bivar}).
  \item \emph{$\icones_{\mathrm{hom}}$}
    (\rocq{Icones.homs.linhom\_functor.linhom\_map\_icones}):
    integral-preservation in $f$
    (\rocq{Icones.homs.linhom\_functor.linhom\_map\_pres\_int}) ---
    $\Lambda(\int \Phi\,d\mu) = \int (\Lambda \circ \Phi)\,d\mu$, by
    uniqueness of the linhom path integral plus $g$'s
    integral-preservation and $D_2$'s Pettis spec.
\end{itemize}
\end{proof}

\begin{definition}[One-sided actions, paper \S5.3]\label{def:linhom-onesided}
\rocq{Icones.homs.linhom\_functor.linhom\_post\_icones}\rocqok
\rocq{Icones.homs.linhom\_functor.linhom\_pre\_icones}\rocqok
\uses{lem:linhom-map-icones}
The covariant post-composition action $C \lin g :=
\texttt{linhom\_map\_icones}\,(\texttt{icones\_id}\,C)\,g : (C \lin
D_1) \to (C \lin D_2)$ is
\rocq{Icones.homs.linhom\_functor.linhom\_post\_icones}, with
$(C \lin g)\,f = g \circ f$
(\rocq{Icones.homs.linhom\_functor.linhom\_post\_iconesE}). The
contravariant pre-composition action $h \lin D :=
\texttt{linhom\_map\_icones}\,h\,(\texttt{icones\_id}\,D) : (C_1 \lin
D) \to (C_2 \lin D)$ is
\rocq{Icones.homs.linhom\_functor.linhom\_pre\_icones}, with
$(h \lin D)\,f = f \circ h$
(\rocq{Icones.homs.linhom\_functor.linhom\_pre\_iconesE}).
\end{definition}

\begin{lemma}[Functoriality of $\lin$, paper Prop 5.8]\label{lem:linhom-functoriality}
\rocq{Icones.homs.linhom\_functor.linhom\_map\_icones\_id}\rocqok
\rocq{Icones.homs.linhom\_functor.linhom\_map\_icones\_comp}\rocqok
\uses{lem:linhom-map-icones}
The hom-functor preserves identities and composition in $\ICones$:
$(\id_C) \lin (\id_D) = \id_{C \lin D}$
(\rocq{Icones.homs.linhom\_functor.linhom\_map\_icones\_id}) and, with
the first slot contravariant,
$(h' \circ h) \lin (g \circ g') = (h' \lin g') \circ (h \lin g)$
(\rocq{Icones.homs.linhom\_functor.linhom\_map\_icones\_comp}). Both
reduce by \texttt{icones\_hom\_eq} to the object-level laws
\texttt{linhom\_map\_id} / \texttt{linhom\_map\_comp} of
Chapter~\ref{ch:linhom}.
\end{lemma}

\section{The tensor product (paper \S5.4)}\label{sec:tensor-build}

Throughout this section we work modulo the staging interface
(\S\ref{sec:tensor-staging}), with local notations $C \lin D :=
\texttt{linhom\_car}\,\Ar\,C\,D$ and $B \otimes C := \texttt{tensor}\,\Ar\,B\,C$.
Everything below is a genuine proof; the only unproved inputs are the
staged \texttt{Parameter}s it consumes.

\subsection{The universal map $\tau$ and the pure tensor}%
\label{sec:tau}

\begin{definition}[The universal map $\tau$ and the pure tensor, paper
\S5.4]\label{def:tau}
\rocq{Icones.homs.tensor.tau}\rocqok
\rocq{Icones.homs.tensor.ptensor}\rocqok
\uses{def:tensor-object, def:linhom-car}
The universal map is $\tau_{B,C} := \Phi_{B,C,B \otimes
C}\bigl(\id_{B \otimes C}\bigr) \in \ICones(B, C \lin B \otimes C)$,
the image of the tensor identity under the hom-bijection $\Phi =
\texttt{tensor\_curry}$ (\rocq{Icones.homs.tensor.tau}). The
\emph{pure tensor} of $x \in B$ and $y \in C$ is
\[
  x \otimes y \;:=\; \tau_{B,C}(x)(y) \;\in\; B \otimes C,
\]
i.e.\ \rocq{Icones.homs.tensor.ptensor}, with the trivial computation
$x \otimes y = \tau_{B,C}(x)(y)$
(\rocq{Icones.homs.tensor.ptensorE}).
\end{definition}

\begin{lemma}[Paper Eq 5.1, $\Phi$ factors through $\tau$]\label{lem:eq51}
\rocq{Icones.homs.tensor.tensor\_curryE}\rocqok
\uses{def:tau}
For $f : \icones_{\mathrm{hom}}\,\Ar\,(B \otimes C)\,D$,
\[
  \Phi(f) \;=\; (C \lin f) \circ \tau_{B,C}
  \qquad\text{(\rocq{Icones.homs.tensor.tensor\_curry\_factor})},
\]
and pointwise $\Phi(f)(x)(y) = f(x \otimes y)$
(\rocq{Icones.homs.tensor.tensor\_curryE}).
\end{lemma}

\begin{proof}
From the concrete naturality of $\Phi$ in $D$
(\rocq{Icones.axioms.saft\_interface.tensor\_curry\_natural\_post},
modulo the staged interface, \S\ref{sec:tensor-staging}) with $h := f$
and $f_0 := \id_{B \otimes C}$: $\Phi(f) = \Phi(f \circ \id) = (C \lin
f) \circ \Phi(\id) = (C \lin f) \circ \tau$. Pointwise, evaluating the
post-composition action through
\rocq{Icones.homs.linhom.linhom\_map\_funE} gives $\Phi(f)(x)(y) =
f(\tau(x)(y)) = f(x \otimes y)$.
\end{proof}

\subsection{Bifunctoriality}\label{sec:tensor-mor}

\begin{definition}[The bifunctor action on morphisms, paper \S5.4]%
\label{def:tensor-mor}
\rocq{Icones.homs.tensor.tensor\_mor}\rocqok
\uses{def:tau, lem:linhom-map-icones}
Given $f : B_1 \to B_2$ and $g : C_1 \to C_2$, the action $f \otimes g
: B_1 \otimes C_1 \to B_2 \otimes C_2$ is the uncurrying of $(C_1 \lin
B_2 \otimes C_2) \xleftarrow{g \lin} (C_2 \lin B_2 \otimes C_2)
\xleftarrow{\tau_{B_2,C_2} \circ f} B_1$, i.e.\
\rocq{Icones.homs.tensor.tensor\_mor}, that precomposes
$\tau_{B_2,C_2} \circ f$ in the hom-slot by $g$.
\end{definition}

\begin{lemma}[Identity law of $\otimes$, paper \S5.4]\label{lem:tensor-mor-id}
\rocq{Icones.homs.tensor.tensor\_mor\_id}\rocqok
\uses{def:tensor-mor}
$\id_B \otimes \id_C = \id_{B \otimes C}$
(\rocq{Icones.homs.tensor.tensor\_mor\_id}), modulo the staged
interface.
\end{lemma}

\begin{proof}
Unfold \texttt{tensor\_mor}; the pre-composition by $\id_C$ collapses
to the identity on $C \lin B \otimes C$ via
$\texttt{linhom\_map\_icones\_id}$
(\autoref{lem:linhom-functoriality}); what remains is
$\texttt{tensor\_uncurry}(\texttt{tensor\_curry}(\id)) = \id$ by the
staged round-trip
\rocq{Icones.axioms.saft\_interface.tensor\_curryK}.
\end{proof}

\begin{remark}[Composition law deferred]
The header of \texttt{theories/homs/tensor.v} delivers only the
\emph{identity} law of the bifunctor; the composition law
$(f_2 \circ f_1) \otimes (g_2 \circ g_1) = (f_2 \otimes g_2) \circ (f_1
\otimes g_1)$ is not mechanised in this iteration. It is not needed
for any coherence diagram of \S\ref{sec:smcc} (each of those is checked
on pure tensors via Prop~5.14), and the SMC-closed bundle
(\autoref{thm:thm515}) records only the identity law in its
\texttt{smcc\_mor\_id} field. The general composition law is deferred
to the M-SAFT discharge together with the rest of the interface.
\end{remark}

\subsection{Theorem 5.12: the hom-tensor adjunction iso}%
\label{sec:thm512}

\begin{theorem}[Paper Theorem 5.12]\label{thm:thm512}
\rocq{Icones.homs.tensor.tensor\_hom\_Phi}\rocqok
\uses{def:icones-iso, def:tau}
For all $B, C, D$, the hom-bijection $\Phi_{B,C,D}$ is an isomorphism
of integrable cones
\[
  (B \otimes C) \lin D \;\cong\; B \lin (C \lin D),
\]
packaged as the \texttt{icones\_iso}
\rocq{Icones.homs.tensor.tensor\_hom\_Phi} (re-exporting the staged
\rocq{Icones.axioms.saft\_interface.tensor\_hom\_iso}, modulo the
interface). Its forward/backward maps are
\rocq{Icones.homs.tensor.tensor\_hom\_fwd} /
\rocq{Icones.homs.tensor.tensor\_hom\_bwd}, and the forward map is
bijective (\rocq{Icones.homs.tensor.tensor\_hom\_fwd\_bij}).
\end{theorem}

At the morphism level, $\Phi = \texttt{tensor\_curry}$ is injective
because \texttt{tensor\_uncurry} is a left inverse: this is
\rocq{Icones.homs.tensor.tensor\_curry\_inj}\rocqok, proved directly
from the staged
\rocq{Icones.axioms.saft\_interface.tensor\_curryK}. Injectivity is the
engine of the extensionality principle below.

\subsection{Theorem 5.13: multiplicativity of the tensor norm}%
\label{sec:thm513}

\begin{theorem}[Paper Theorem 5.13]\label{thm:thm513}
\rocq{Icones.homs.tensor.tensor\_normME}\rocqok
\uses{def:tau}
For $x \in B$ and $y \in C$,
\[
  \cnorm{x \otimes y} \;=\; \cnorm{x} \cdot \cnorm{y}.
\]
\end{theorem}

\begin{proof}[Proof sketch]
The $\le$ direction is genuinely derived
(\rocq{Icones.homs.tensor.tensor\_norm\_le}) from norm-decrease of
$\tau$ (an $\icones_{\mathrm{hom}}$, hence a $\cones_{\mathrm{hom}}$):
$\cnorm{x \otimes y} = \cnorm{\tau(x)(y)} \le \cnorm{\tau(x)} \cdot
\cnorm{y} \le \cnorm{x} \cdot \cnorm{y}$. The converse $\ge$ needs the
dual-separation Prop~3.11 together with the element-level inverse
$\Phi^{-1}$ of \autoref{thm:thm512}, both supplied by the SAFT
discharge; the full equality
\rocq{Icones.homs.tensor.tensor\_normME} therefore rests on the staged
\rocq{Icones.axioms.saft\_interface.tensor\_normM} (modulo the
interface, \S\ref{sec:tensor-staging}).
\end{proof}

\subsection{Proposition 5.14: pure-tensor extensionality}%
\label{sec:tensor-ext}

This is the lever that makes all monoidal coherence \emph{provable}
from the staged computation laws: a morphism out of an iterated tensor
is determined by its values on (iterated) pure tensors.

\begin{lemma}[Paper Proposition 5.14, binary case]\label{lem:tensor-ext}
\rocq{Icones.homs.tensor.tensor\_ext}\rocqok
\uses{thm:thm512, lem:eq51}
Two morphisms $f, g : B \otimes C \to D$ that agree on every pure
tensor $x \otimes y$ are equal.
\end{lemma}

\begin{proof}
By injectivity of $\Phi$
(\rocq{Icones.homs.tensor.tensor\_curry\_inj}) it suffices that
$\Phi(f) = \Phi(g)$; by \texttt{icones\_hom\_eq} and
\texttt{linhom\_eq} this is pointwise; and by Eq~5.1
(\autoref{lem:eq51}) $\Phi(f)(x)(y) = f(x \otimes y) = g(x \otimes y) =
\Phi(g)(x)(y)$.
\end{proof}

\begin{lemma}[Paper Proposition 5.14, ternary tree]\label{lem:tensor-ext3}
\rocq{Icones.homs.tensor.tensor\_ext3}\rocqok
\uses{lem:tensor-ext}
Two morphisms $f, g : (A \otimes B) \otimes C \to D$ that agree on
every iterated pure tensor $(x \otimes y) \otimes z$ are equal.
\end{lemma}

\begin{proof}
The paper's induction at depth two: by injectivity of $\Phi$ it
suffices that $\Phi(f) = \Phi(g) : A \otimes B \to (C \lin D)$; by the
binary \autoref{lem:tensor-ext} it suffices they agree on $x \otimes
y$; by \texttt{linhom\_eq} on every $z : C$; and by Eq~5.1 $\Phi(f)(x
\otimes y)(z) = f\bigl((x \otimes y) \otimes z\bigr)$, equal to the
$g$-side by hypothesis.
\end{proof}

A quaternary instance
\rocq{Icones.homs.smcc.tensor\_ext4}\rocqok (one more currying layer
than \autoref{lem:tensor-ext3}, for the tree $((A \otimes B) \otimes C)
\otimes D$) is proved the same way and is what the pentagon needs.

\section{The symmetric monoidal closed structure (paper \S5.5)}%
\label{sec:smcc}

We assemble the SMC closed structure on top of $\otimes$ and the staged
structural isos, with $1 := \texttt{cone\_one\_car}\,\Ar$ the monoidal
unit (the integrable cone $\Rnn$ of Chapter~\ref{ch:icones}).

\subsection{Bilinearity of the pure tensor}\label{sec:tensor-bilin}

\begin{lemma}[Bilinearity of $\otimes$, paper \S5.4]\label{lem:tensor-bilin}
\rocq{Icones.homs.smcc.ptensorZl}\rocqok
\rocq{Icones.homs.smcc.ptensorZr}\rocqok
\uses{def:tau}
The pure tensor is linear in each argument. A nonnegative scalar
commutes through either slot:
$(r \cdot x) \otimes y = r \cdot (x \otimes y)$
(\rocq{Icones.homs.smcc.ptensorZl}) and $x \otimes (r \cdot y) = r
\cdot (x \otimes y)$ (\rocq{Icones.homs.smcc.ptensorZr}). Linearity in
$y$ is linearity of the linhom $\tau(x)$; linearity in $x$ uses that
$\tau$ is itself a linear morphism (it scales $\tau(x)$ pointwise, via
\rocq{Icones.homs.smcc.linhom\_funZ}).
\end{lemma}

\subsection{The structural isos and their computation laws}%
\label{sec:struct-isos}

\begin{definition}[The structural isos $\alpha/\lambda/\rho/\sigma$,
paper Eqs 5.2--5.4]\label{def:struct-isos}
\rocq{Icones.homs.smcc.tensor\_assoc}\rocqok
\rocq{Icones.homs.smcc.tensor\_lunit}\rocqok
\rocq{Icones.homs.smcc.tensor\_runit}\rocqok
\rocq{Icones.homs.smcc.tensor\_braid}\rocqok
\uses{def:icones-iso, sec:tensor-staging}
The associator $\alpha_{A,B,C} : (A \otimes B) \otimes C \cong A
\otimes (B \otimes C)$
(\rocq{Icones.homs.smcc.tensor\_assoc}), the left/right unitors
$\lambda_A : 1 \otimes A \cong A$
(\rocq{Icones.homs.smcc.tensor\_lunit}) and $\rho_A : A \otimes 1 \cong
A$ (\rocq{Icones.homs.smcc.tensor\_runit}), and the braiding
$\sigma_{A,B} : A \otimes B \cong B \otimes A$
(\rocq{Icones.homs.smcc.tensor\_braid}) are the
\texttt{icones\_iso}s re-exported under their paper names from the
staged structural-iso \texttt{Parameter}s. Each comes with its
pure-tensor computation law:
\[
  \alpha\bigl((x \otimes y) \otimes z\bigr) = x \otimes (y \otimes z),
  \quad
  \lambda(u \otimes x) = u \cdot x = \rho(x \otimes u),
  \quad
  \sigma(x \otimes y) = y \otimes x,
\]
namely \rocq{Icones.homs.smcc.tensor\_assocEp},
\rocq{Icones.homs.smcc.tensor\_lunitEp},
\rocq{Icones.homs.smcc.tensor\_runitEp},
\rocq{Icones.homs.smcc.tensor\_braidEp} (here $u \in 1 = \Rnn$ acts by
scaling). These restate the staged computation laws
\rocq{Icones.axioms.saft\_interface.tensor\_assocE} etc.\ (modulo the
interface, \S\ref{sec:tensor-staging}).
\end{definition}

The bifunctor action computes on pure tensors:
$(f \otimes g)(x \otimes y) = (f\,x) \otimes (g\,y)$, namely
\rocq{Icones.homs.smcc.tensor\_morE}\rocqok, derived from the staged
round-trip \rocq{Icones.axioms.saft\_interface.tensor\_uncurryK} and
Eq~5.1.

\subsection{The derived coherence laws}\label{sec:coherence}

The point of \S\ref{sec:tensor-ext} is that, although the structural
\emph{isos} are staged, every coherence \emph{law} between them is
genuinely \emph{proved}, by checking the two sides agree on (iterated)
pure tensors and reading off the computation laws.

\begin{lemma}[Symmetry law, paper Eq 5.4]\label{lem:braid-invol}
\rocq{Icones.homs.smcc.tensor\_braid\_invol}\rocqok
\uses{def:struct-isos, lem:tensor-ext}
$\sigma_{B,A} \circ \sigma_{A,B} = \id_{A \otimes B}$.
\end{lemma}
\begin{proof}
By \autoref{lem:tensor-ext}: $\sigma(\sigma(x \otimes y)) = \sigma(y
\otimes x) = x \otimes y$ via \texttt{tensor\_braidEp}.
\end{proof}

\begin{lemma}[Triangle coherence, paper Thm 5.15]\label{lem:triangle}
\rocq{Icones.homs.smcc.tensor\_triangle}\rocqok
\uses{def:struct-isos, lem:tensor-bilin, lem:tensor-ext3}
$(\id_A \otimes \lambda_B) \circ \alpha_{A,1,B} = \rho_A \otimes \id_B$
as maps $(A \otimes 1) \otimes B \to A \otimes B$.
\end{lemma}
\begin{proof}
By \autoref{lem:tensor-ext3}: on $(x \otimes u) \otimes y$ both sides
equal $u \cdot (x \otimes y)$ --- the LHS through the right slot
(\texttt{tensor\_assocEp}, \texttt{tensor\_morE},
\texttt{tensor\_lunitEp}, \texttt{ptensorZr}), the RHS through the left
slot (\texttt{tensor\_morE}, \texttt{tensor\_runitEp},
\texttt{ptensorZl}).
\end{proof}

\begin{lemma}[Pentagon coherence, paper Thm 5.15]\label{lem:pentagon}
\rocq{Icones.homs.smcc.tensor\_pentagon}\rocqok
\uses{def:struct-isos, lem:tensor-ext3}
The pentagon identity for $\alpha$, as maps $((A \otimes B) \otimes C)
\otimes D \to A \otimes (B \otimes (C \otimes D))$.
\end{lemma}
\begin{proof}
By the quaternary \texttt{tensor\_ext4}: on $((w \otimes x) \otimes y)
\otimes z$ both composites equal $w \otimes (x \otimes (y \otimes z))$,
by repeated \texttt{tensor\_assocEp} and the bifunctor law
\texttt{tensor\_morE}.
\end{proof}

\begin{lemma}[Braiding hexagon, paper Thm 5.15]\label{lem:hexagon}
\rocq{Icones.homs.smcc.tensor\_hexagon}\rocqok
\uses{def:struct-isos, lem:tensor-ext3}
The hexagon identity for $\sigma$, as maps $(A \otimes B) \otimes C \to
B \otimes (C \otimes A)$.
\end{lemma}
\begin{proof}
By \autoref{lem:tensor-ext3}: on $(x \otimes y) \otimes z$ both
composites equal $y \otimes (z \otimes x)$, by \texttt{tensor\_assocEp},
\texttt{tensor\_braidEp} and \texttt{tensor\_morE}.
\end{proof}

\subsection{Theorem 5.15: $\ICones$ is symmetric monoidal closed}%
\label{sec:thm515}

\begin{theorem}[Paper Theorem 5.15]\label{thm:thm515}
\rocq{Icones.homs.smcc.ICones\_smcc}\rocqok
\uses{def:tensor-object, lem:tensor-mor-id, def:struct-isos,
lem:braid-invol, lem:triangle, lem:pentagon, lem:hexagon, thm:thm512}
$\ICones$ is a symmetric monoidal closed category. The structure is
bundled in the record
\rocq{Icones.homs.smcc.ICones\_SMCC}, whose canonical witness
\rocq{Icones.homs.smcc.ICones\_smcc} populates every field with the
proved data:
\begin{itemize}
  \item the monoidal product $\otimes$ (\texttt{tensor}) with bifunctor
    action $\otimes$ on morphisms (\texttt{tensor\_mor}) and its
    identity law (\texttt{tensor\_mor\_id}), and the unit $1 =
    \texttt{cone\_one\_car}\,\Ar$;
  \item the structural isos $\alpha/\lambda/\rho/\sigma$
    (\autoref{def:struct-isos});
  \item the symmetry law $\sigma \circ \sigma = \id$
    (\autoref{lem:braid-invol}) and the triangle, pentagon and
    braiding-hexagon coherence diagrams
    (\autoref{lem:triangle}, \autoref{lem:pentagon},
    \autoref{lem:hexagon});
  \item closedness: the internal hom $\lin$ ($\texttt{linhom\_car}$)
    with the Theorem~5.12 iso $(B \otimes C) \lin D \cong B \lin (C
    \lin D)$ (\autoref{thm:thm512}).
\end{itemize}
\end{theorem}

\begin{proof}
Each field is one of the proved lemmas above; the bundle
\rocq{Icones.homs.smcc.ICones\_smcc} just assembles them. The only
unproved inputs are the staged structural isos and their computation
laws (paper Eqs~5.2--5.4) and the universal property of $\otimes$,
discharged with the rest of the SAFT contract.
\end{proof}

\section{Status}\label{sec:tensor-status}

\paragraph{Delivered (mechanised modulo the interface).}
The three prerequisites of \S5.3 are fully proved with no project
axioms: isos in $\ICones$ with their groupoid structure
(\rocq{Icones.homs.icones\_iso.icones\_iso}), the Def~3.13 test-pullback
characterisation
(\rocq{Icones.mcones.test\_pullback.test\_pullback\_meas},
\rocq{Icones.mcones.test\_pullback.test\_pullback\_meas\_bivar},
\rocq{Icones.mcones.test\_pullback.test\_pullback\_cont}), and the
morphism-level internal-hom functor of Prop~5.8
(\rocq{Icones.homs.linhom\_functor.linhom\_map\_icones}, its one-sided
actions, and functoriality). On top of the staging interface, the
tensor theory of \S5.4 (the universal map $\tau$ and pure tensor, Eq~5.1,
bifunctor identity law, Thm~5.12, the $\le$ half of Thm~5.13, and the
binary/ternary/quaternary pure-tensor extensionality of Prop~5.14) and
the SMC closed structure of \S5.5 (the structural isos, the
\emph{proved} symmetry/triangle/pentagon/hexagon coherence, and the
bundle \rocq{Icones.homs.smcc.ICones\_smcc} witnessing Theorem~5.15)
are all mechanised --- with the substantive coherence content genuinely
\Qed{}'d, not assumed.

\paragraph{Staged (to be discharged by M-SAFT).}
The deep representability content remains the $\approx 18$
\texttt{Parameter}s of \texttt{theories/axioms/saft\_interface.v}
(\S\ref{sec:tensor-staging}): the tensor object former
(\rocq{Icones.axioms.saft\_interface.tensor}), the hom-bijection and its
naturality (\rocq{Icones.axioms.saft\_interface.tensor\_curry},
\rocq{Icones.axioms.saft\_interface.tensor\_uncurry},
\rocq{Icones.axioms.saft\_interface.tensor\_curry\_natural\_post}),
the Thm~5.12 iso (\rocq{Icones.axioms.saft\_interface.tensor\_hom\_iso}),
the Thm~5.13 norm equation
(\rocq{Icones.axioms.saft\_interface.tensor\_normM}), and the structural
isos of Eqs~5.2--5.4
(\rocq{Icones.axioms.saft\_interface.tensor\_assoc\_iso},
\rocq{Icones.axioms.saft\_interface.tensor\_lunit\_iso},
\rocq{Icones.axioms.saft\_interface.tensor\_runit\_iso},
\rocq{Icones.axioms.saft\_interface.tensor\_braid\_iso}) with their
computation laws. The plan
(\texttt{PLAN.md} \S13.1, ``M-SAFT'') is to discharge every
\texttt{Parameter} by proving the Special Adjoint Functor Theorem for
$\ICones$, after which \texttt{saft\_interface.v} is deleted and the
project returns to zero project axioms. The bifunctor composition law
(\S\ref{sec:tensor-mor}) and the naturality of $\Phi$ in the $C$ slot
(\S\ref{sec:tensor-staging}) are deferred together with the discharge.
The MVP headline theorem (paper Theorem~6.5,
Chapter~\ref{ch:skern}) is independent of this interface.
